% from main.tex
%%%%%%%%%%  scan idx 492  |  folio 479  %%%%%%%%%%
factorisation de $f$ comme composé d'une immersion et d'un morphisme lisse, le complexe $L^{X/Y}_\bullet$ défini en 2.1 est strictement parfait.

Soit $f=g\circ i$, où $g$ est un morphisme lisse, et $i$ une immersion. D'après 1.2, $i$ est une immersion régulière. Si $J$ est l'Idéal définissant $i$, $J/J^2$ est donc un $\mathcal O_X$-Module localement libre de type fini. De même, comme $g:V\to Y$ est lisse, $i^*(\Omega^1_{V/Y})$ est localement libre de type fini, d'où 2.4.

\textbf{Corollaire 2.5.} Soit $f:X\to Y$ un morphisme d'intersection complète lissifiable, et $f=g\circ i$ une factorisation de $f$ par un $Y$-schéma lisse $V$. Soit $J$ l'idéal de $\mathcal O_V$ définissant $i$. Alors, si $K^{\bullet}(X)$ désigne le groupe de Grothendieck des $\mathcal O_X$-Modules localement libres de type fini (resp. des complexes parfaits sur $X$), l'élément $T_f\in K^{\bullet}(X)$ défini par :
\[
T_f=\operatorname{cl}(i^*(\Omega^1_{V/Y}))-\operatorname{cl}(J/J^2)
\]
(resp.
\[
T_f=\operatorname{cl}(L^{X/Y}_\bullet)
\]
)
ne dépend pas de la factorisation de $f$ choisie. La valeur en un point $x$ de $X$ de l'augmentation de $T_f$ n'est autre que la dimension relative virtuelle de $X$ sur $Y$ au point $x$ (1.9).

L'élément $T_f$ est bien défini dans les deux cas envisagés d'après 2.4, et sa classe dans $K^{\bullet}(X)$ ne dépend pas de la factorisation, d'après 2.2, dans le cas respé, et d'après la démonstration de 2.2 dans le cas non respé. La dernière assertion de 2.5 est triviale.

\textbf{Proposition 2.6.} Soient $f:X\to Y$, et $g:Y\to Z$ deux morphismes d'intersection complète lissifiables, et supposons de plus la condition suivante remplie:

